\begin{frame}{왜 합성 데이터셋인가 + 레이어링의 특징}
  \begin{columns}[T]
    \begin{column}{0.5\textwidth}
      \begin{block}{합성인 것의 의미}
        \small ``\textbf{답이 알려진} 거래 그래프''를 만들어 탐지
        모델의 성능을 \textbf{측정}할 수 있게 하는 것 --- 실제
        금융 데이터는 \textbf{프라이버시} 때문에 공개가 어려우니,
        \textbf{구조와 노이즈 수준만 조절}해 ``탐지 가능한 패턴''을
        심어놓은 합성 데이터셋.
      \end{block}
      \vskip0.5em
      \begin{itemize}
        \item 레이어링 패턴의 특징: 불법 자금의 흐름이
          \textbf{1--2홉이 아니라 3--4홉 이상}으로 뻗어 있다.
      \end{itemize}
    \end{column}
    \begin{column}{0.47\textwidth}
      \centering
      \includegraphics[width=\linewidth]{aml_layering_graph.png}
    \end{column}
  \end{columns}

\note{\begin{itemize}
\item ``계좌 A가 B로, B가 C로, C가 해외로'' 이동하는 동안
      \textbf{각 구간은 정상 거래와 구분하기 어렵다}.
\end{itemize}}
\end{frame}
